\section{Eager Normal-Form Bisimulation}
\label{s:enf}


From Corollary~\ref{c:OGS_COGS} and Theorem~\ref{t:adrien},
we can immediately conclude that the semantics on $\lambda$-terms induced by OGS
(Alternating or Concurrent) coincides with 
enf-bisimilarity (i.e., Lassen's trees), as introduced in Definition~\ref{d:enf}. 

In this section, we show a direct proof of the result, for \COGS bisimilarity, 
as an example of application of the  
up-to composition technique for \COGS
 presented in Section~\ref{ss:upto_games}.

% We prove that enf-bisimulations, as presented in Section~\ref{s:back:trees},
% can be lifted to a bisimulation up-to composition over the Concurrent OGS LTS,
% as presented in Section~\ref{ss:upto_games}.


Terms, values, and evaluation contexts  of the $\lambda$-calculus can be directly lifted to singleton concurrent configurations,
meaning that we can transform a relation on terms, values, and contexts $\R$ into
a relation $\liftRel{\R}$ on singleton concurrent configurations in the following way:
\begin{itemize}
 \item If $(M_1,M_2) \in \R$ then $(\conf{\pmap{p}{M_1},\emptymap,\phi},
 \conf{\pmap{p}{M_2},\emptymap,\phi}) \in \liftRel{\R}$, 
 with $\phi = \fv{M_1,M_2} \uplus \{p\}$;
%
 \item If $(V_1,V_2) \in \R$ then $(\conf{\pmap{x}{V_1},\phi},\conf{\pmap{x}{V_2},\phi})
 \in \liftRel{\R}$, with $\phi= \fv{V_1,V_2} \uplus \{x\}$;
%
 \item If $(E_1,E_2) \in \R$ then $(\conf{\pmap{p}{(E_1,q)},\phi},
  \conf{\pmap{p}{(E_2,q)},\phi}) \in \liftRel{\R}$, with 
  $\phi= \fv{E_1,E_2} \uplus \{p,q\}$.
\end{itemize}

\begin{theorem}
 Taking $(\RR_{\enfM},\RR_{\enfV},\RR_{\enfK})$ an enf-bisimulation, then 
 $(\liftRel{\RR_{\enfM}} \cup \liftRel{\RR_{\enfV}} \cup \liftRel{\RR_{\enfK}})$
 is a singleton bisimulation up-to composition in \COGS (as defined in Section~\ref{ss:upto_games}).
\end{theorem}

\begin{proof}
 Taking $\FF_1,\FF_2$ two singleton configurations s.t. 
 $(\FF_1,\FF_2) \in \liftRel{\RR_{\enfM}}$,
 then we can write $\FF_i$ as $\conf{\pmap{p}{M_i},\phi}$
 for $i=1,2$, such that $(M_1,M_2) \in \R$.
 Suppose that $\FF_1 \Longrightarrowp\arrC\ell 
%\ArrC{\act}
 \FF'_1 $,
 with $\ell$ a Player action. We only present the Player Question case,
which is where  `up-to composition' is useful.
So writing $\ell$ as $\questP{x}{y,q}$,
$\FF'_1 $ can be written as $\conf{[y \mapsto V_1]\cdot[q \mapsto (E_1,p)]}$,
so that $M_1 \redv^{*} E_1[x V_1]$.
 
As $(M_1,M_2) \in \R$, there are $E_2,V_2$ s.t.
 $M_2 \redv^{*} E_2[x V_2]$, $(V_1,V_2) \in \R$ and $(E_1,E_2) \in \R$.
Hence $\FF_2    
\Longrightarrowp\arrC\ell
% \ArrC{\act}
 \FF'_2$ with $\FF'_2 = 
 \conf{[y \mapsto V_2]\cdot[q \mapsto (E_2,p)]}$.
 
Finally, $(\conf{\pmap{y}{V_1}},\conf{\pmap{y}{V_2}}) \in \liftRel{\R}$ 
 and $(\conf{\pmap{q}{(E_1,p)}},\conf{\pmap{q}{(E_2,p)}}) \in \liftRel{\R}$,
 so that 
 ${(\FF'_1,\FF'_2)} \in {\uptoComp{\liftRel{\R}}}$.
 
 The case of passive singleton configurations is proved in a similar way.
 \end{proof}
